\section{Реализуемость: лестница импеданса}
\label{sec:realizability}

Теорема субстратной замкнутости УГМ (T-153) доказывает, что логическая
архитектура работает на любом субстрате, удовлетворяющем трём условиям --- (C1)
конечномерное, декогеренция-свободное эффективное пространство состояний; (C2)
CPTP-(линдбладова) динамика; (C3) не менее семи некоммутирующих наблюдаемых мод.
Поэтому нет единственной машины; есть \emph{семейство}, и его члены различаются
лишь тем, насколько хорошо их физические примитивы соответствуют логической
структуре. Мы ранжируем их по \emph{импедансу} --- рассогласованию, которое
субстрат должен замазать --- против четырёх критериев соответствия:

\begin{itemize}[leftmargin=1.5em,itemsep=0pt]
  \item[\textbf{(M1)}] \emph{эволюция на месте} --- состояние продвигается без
        цикла выборка--исполнение;
  \item[\textbf{(M2)}] \emph{Фано-проводка интерконнекта} --- семь тройственных
        шин, не кроссбар;
  \item[\textbf{(M3)}] \emph{диссипативно--регенеративная родная динамика} ---
        физика уже есть Линдблад$+$гомеостаз;
  \item[\textbf{(M4)}] \emph{семь некоммутирующих мод} (C3).
\end{itemize}

\begin{table}[H]
\centering\small
\caption{Лестница реализуемости. Высшие ступени соответствуют логической
структуре с меньшим импедансом (меньше слоёв эмуляции); низшие ступени строимы
сегодня. Заявление \status{С}: теорема фиксирует логическую архитектуру и условия
C1--C3, не победивший субстрат.}
\label{tab:ladder}
\begin{tabular}{@{}p{4.3cm}cccc p{3.0cm}@{}}
\toprule
\textbf{Субстрат} & \textbf{M1} & \textbf{M2} & \textbf{M3} & \textbf{M4} & \textbf{Статус} \\
\midrule
Управляемо-диссипативное аналог./фотонное/открыто-квантовое 7-модовое устройство & \checkmark & \checkmark & \checkmark & \checkmark & наименьший импеданс; исследование \\
Мемристорный Фано-кроссбар (аналог.\ MAC в памяти на 7 шинах) & \checkmark & \checkmark & частично & \checkmark & ближнесрочно \\
Нейроморфное спайковое ядро, Фано-топология, событийное & \checkmark & \checkmark & частично & \checkmark & ближнесрочно \\
GPU / CPU с \DM{} в L1 и Фано-разреженными ядрами & эмул. & эмул. & эмул. & \checkmark & \textbf{существует (\SYNARC{})} \\
\bottomrule
\end{tabular}
\end{table}

\begin{figure}[H]
\centering
\begin{tikzpicture}[font=\scriptsize, >=Stealth]
  \draw[very thick, talossteel] (0,0) -- (1.4,4.9);
  \draw[very thick, talossteel] (3.2,0) -- (4.6,4.9);
  \foreach \t in {0.06,0.36,0.66,0.96}{
    \draw[thick, talosamber] ({0+1.4*\t},{4.9*\t}) -- ({3.2+1.4*\t},{4.9*\t}); }
  \node[anchor=west] at (4.9,0.28)  {\textbf{Ступень 0 --- эмуляция GPU/CPU (существует).} корректность + телеметрия; данные всё же движутся};
  \node[anchor=west] at (5.3,1.75)  {\textbf{Ступень 1 --- мемристорный Фано-кроссбар.} обновление в массиве; первый честный тест B1};
  \node[anchor=west] at (5.7,3.22)  {\textbf{Ступень 2 --- нейроморфное Фано-ядро.} событийное, на месте};
  \node[anchor=west] at (6.1,4.68)  {\textbf{Ступень 3 --- управляемо-диссипативный прибор.} физика $=$ тик; у энергопола};
  \draw[->, thick, plosgray] (-0.7,0.1) -- node[left, align=center]{\tiny импеданс\\ \tiny падает} (-0.1,4.7);
\end{tikzpicture}
\caption{Лестница реализуемости как лестница. Каждая ступень исполняет
тождественную логическую архитектуру (T-153) и тот же софт; подъём убирает
слои эмуляции (импеданс падает, стрелка слева) и усиливает энергетическое
заявление ступень за ступенью. Номера ступеней совпадают с дорожной картой
\S\ref{sec:discussion}.}
\label{fig:ladderfig}
\end{figure}

Нижняя ступень не гипотетична: это текущая реализация \SYNARC{}, держащая
\DM{} ($784$~байта) в L1-кэше, продвигающая её Странг-расщеплённым
Фано-разреженным ядром при $\sim\!700$--$1400$~FLOP на тик и арифметической
глубине $\log_2 7\approx 2{,}8$, и воспроизводящая всякий инвариант и точный
закон бит-в-бит против независимого оракула. Она эмулирует M1--M3 в софте
(наибольший импеданс) и поэтому не реализует энергетический пол \S\ref{sec:energy}
--- GPU всё же двигает данные --- но доказывает, что логическая архитектура верна
и запускаема сегодня. Каждая высшая ступень убирает слой эмуляции: мемристорный
Фано-кроссбар выполняет обновление когерентности \emph{в} массиве (M1, M2 родные);
управляемо-диссипативное аналоговое устройство имеет линдбладову динамику как
\emph{родную} физику (M3 родная) и приближается к энергетическому полу, потому
что его обратимое ядро буквально обратимо. Лестница --- дорожная карта, не меню:
по ней взбираются по мере того, как позволяет производство, а софт (\SYNARC{}) не
меняется на каждой ступени, потому что логическая архитектура фиксированна (T-153).
